\section{Proposed Methodology}

\subsection{Overall Framework}
Quy trình nghiên cứu đề xuất gồm sáu bước chính. Thứ nhất, dữ liệu S3 được chuẩn hóa hình ảnh và loại bỏ các mẫu không có nhãn định danh cụ thể ở cấp loài. Thứ hai, các hình ảnh được nhóm theo thư mục con (subfolder) đại diện cho mẫu vật lý. Thứ ba, các đặc trưng biểu diễn ban đầu được trích xuất bằng mạng trích xuất đặc trưng cơ sở (backbone) mạnh như EfficientNetV2-M hoặc Swin-Large để lập bản đồ phân bố ngữ nghĩa của từng mẫu vật. Thứ tư, giao thức CEGS-Split tự động chọn phương pháp phân chia dữ liệu tối ưu cho từng loài gỗ nhằm thiết lập các tập huấn luyện (train), kiểm định (validation) và kiểm thử (test) cô lập mẫu vật lý, đảm bảo chống rò rỉ dữ liệu. Thứ năm, mạng ConvNeXt-Tiny được tinh chỉnh (fine-tune) kết hợp với bộ đầu chiếu (projection head) phi tuyến để học biểu diễn đặc trưng nhúng 256 chiều. Thứ sáu, mô hình được đánh giá toàn diện bằng các phép truy xuất hình ảnh (retrieval), thuật toán phân cụm (clustering), trực quan hóa giảm chiều t-SNE và bản đồ kích hoạt chú ý Grad-CAM/Finer-CAM.

\subsection{Embedding-Aware Class-wise Specimen Split (CEGS-Split)}
Với mỗi thư mục con $S_j$, đặc trưng nhúng đại diện được tính bằng trọng tâm (centroid) của các hình ảnh thuộc thư mục con đó:
\begin{equation}
E_j=\frac{1}{|S_j|}\sum_{x_i\in S_j} e_i,
\end{equation}
trong đó $e_i\in\mathbb{R}^{D}$ là vector đặc trưng của hình ảnh $x_i$. Nếu số chiều quá lớn so với số lượng thư mục con, thuật toán phân tích thành phần chính (PCA) được áp dụng để giảm chiều về $d'=\min(G-2,128)$ nhằm ổn định việc ước lượng ma trận hiệp phương sai và giảm thiểu hiện tượng quá khớp do số chiều lớn (curse of dimensionality).

Giao thức CEGS-Split hoạt động bằng cách quét qua một không gian các chiến lược phân chia tiềm năng (gồm các nhóm phương pháp chính: Mahalanobis distance, Hierarchical Clustering, phân nhóm có phân tầng và phân cụm phân tầng) để xác định thuật toán phù hợp nhất cho từng loài gỗ cụ thể. Trong đó, bốn chiến lược phân chia cốt lõi được phát triển nhằm duy trì tính cô lập thư mục con nghiêm ngặt ở các mức độ hình học và phân phối khác nhau. Phép phân chia ngẫu nhiên có phân tầng đóng vai trò phương pháp đối chiếu cơ sở (baseline) để so sánh hiệu năng, còn các phương pháp còn lại đều cố gắng giữ tính cô lập mẫu vật lý ở các mức độ khác nhau nhằm triệt tiêu rò rỉ thông tin. Phương pháp CEGS-Split tự động chọn phương pháp phân chia tối ưu theo từng loài thay vì áp dụng một quy tắc cứng nhắc duy nhất cho toàn bộ bộ dữ liệu.

\subsubsection{Mahalanobis Distance-based Split}
Chiến lược này xem các mẫu vật ở xa trọng tâm (centroid) của lớp là các mẫu dị biệt (outlier) về mặt ngữ nghĩa và đưa chúng vào tập kiểm thử hoặc kiểm định nhằm thử nghiệm giới hạn (stress-test) khả năng tổng quát hóa của mô hình. Với tập mẫu chưa phân bổ $U$, centroid $\mu_U$ và hiệp phương sai $\Sigma_U$ được xác định bởi:
\begin{equation}
\mu_U=\frac{1}{|U|}\sum_{j\in U} z_j,
\end{equation}
\begin{equation}
\Sigma_U=\frac{1}{|U|-1}\sum_{j\in U}(z_j-\mu_U)(z_j-\mu_U)^T+\epsilon I.
\end{equation}
Khoảng cách Mahalanobis của mẫu $j$ đến centroid hiện tại là:
\begin{equation}
d_{\text{Mahal}}(z_j,\mu_U)=\sqrt{(z_j-\mu_U)^T\Sigma_U^{-1}(z_j-\mu_U)}.
\end{equation}
Tại mỗi vòng lặp, mẫu ở xa nhất $j^*=\arg\max_{j\in U}d_{\text{Mahal}}(z_j,\mu_U)$ được đưa vào tập kiểm thử hoặc kiểm định dựa trên kích thước yêu cầu, các mẫu còn lại được giữ lại làm dữ liệu huấn luyện.

\subsubsection{Hierarchical Clustering-based Split}
Chiến lược này phát hiện các nhóm mẫu vật tự nhiên trong cùng một loài bằng thuật toán Hierarchical Clustering liên kết tích tụ (Agglomerative Hierarchical Clustering) với tiêu chí liên kết Ward. Số lượng cụm tối ưu $K^*$ được xác định bằng phương pháp khuỷu tay (elbow method) áp dụng trên tổng bình phương khoảng cách trong cụm (Within-Cluster Sum of Squares - WCSS):
\begin{equation}
\text{WCSS}(K)=\sum_{k=1}^{K}\sum_{z_j\in\mathcal{C}_k}\|z_j-\mu_{\mathcal{C}_k}\|_2^2.
\end{equation}
After scaling $K$ and $\text{WCSS}(K)$ to $[0,1]$, the elbow point is selected by:
\begin{equation}
K^*=\arg\max_K \frac{|x_K+y_K-1|}{\sqrt{2}}.
\end{equation}
Mỗi cụm $\mathcal{C}_k$ được giữ nguyên vẹn khi phân bổ vào tập huấn luyện, kiểm định hoặc kiểm thử. Điều này ngăn các nhóm vân gỗ có cấu trúc bề mặt gần nhau bị chia cắt sang các tập khác nhau, giảm thiểu nguy cơ rò rỉ dữ liệu ở cấp độ cụm (cluster-level leakage).

\subsubsection{Stratified Group-based Split}
Phương pháp phân chia nhóm có phân tầng (Stratified Group Split) tối ưu hóa đồng thời hai ràng buộc: cô lập nhóm mẫu vật và cân bằng phân bố lớp giữa các tập. Gọi $w_{g,c}$ là số lượng ảnh của loài $c$ trong nhóm $g$, số lượng mẫu mục tiêu cho phân chia $f$ là:
\begin{equation}
T_{f,c}=\frac{1}{K}\sum_{g=1}^{G}w_{g,c}.
\end{equation}
Thuật toán gán tham lam sẽ lựa chọn tập:
\begin{equation}
f^*=\arg\min_f \sum_{c=1}^{C}\left(\frac{C_{f,c}+w_{g,c}}{T_{f,c}}-1\right)^2,
\end{equation}
trong đó $C_{f,c}$ là số lượng mẫu hiện tại của loài $c$ trong tập $f$. Phương pháp phân chia này phù hợp với các loài gỗ cần ưu tiên sự cân bằng số lượng mẫu giữa các tập để ổn định quá trình tính toán đạo hàm (gradient).

\subsubsection{Agglomerative Stratified Split}
Phương pháp này kết hợp giữa phân cụm đặc trưng và phân chia phân tầng theo khoảng cách hình học. Sau khi gom cụm mẫu vật, khoảng cách từ trọng tâm cụm $\mu_{\mathcal{C}_k}$ đến trọng tâm chung của toàn loài $\mu_{\text{global}}$ được tính bằng:
\begin{equation}
d(\mathcal{C}_k) = \sqrt{ D^2(\mu_{\mathcal{C}_k}, \mu_{\text{global}}) },
\end{equation}
where:
\begin{equation}
D^2(\mu_{\mathcal{C}_k}, \mu_{\text{global}}) = (\mu_{\mathcal{C}_k}-\mu_{\text{global}})^T\Sigma_{\text{global}}^{-1}(\mu_{\mathcal{C}_k}-\mu_{\text{global}}).
\end{equation}
Các cụm được chia thành ba dải Near, Mid và Far. Khi phân bổ, thuật toán chọn tập có độ hụt lớn nhất so với mục tiêu:
\begin{equation}
\Delta_s=\max(0,N_s^{\text{target}}-N_s^{\text{current}}),\quad s\in\{\text{train},\text{val},\text{test}\}.
\end{equation}
Nhờ đó, cả ba tập huấn luyện, kiểm định và kiểm thử đều chứa đầy đủ các mẫu điển hình, mẫu trung gian và mẫu dị biệt, tránh việc dồn tất cả các mẫu ngoại lai (outlier) vào tập kiểm thử.

\subsection{CEGS-Split Configuration}
CEGS-Split là một cấu hình lai theo từng loài. Thay vì áp dụng một chiến lược duy nhất cho toàn bộ dataset, mỗi loài được gán phương pháp chia phù hợp với đặc điểm sinh học, số lượng subfolder và kết quả phân tích phân bố không gian biểu diễn. Cấu hình này được suy ra từ quy trình tìm kiếm tự động quét các chiến lược chia dữ liệu trên từng lớp, dùng EfficientNetV2-M để tạo embedding phục vụ thuật toán phân hoạch, dùng Swin-Large để đo nearest-neighbor similarity giữa train và test, rồi chọn phương pháp có mean nearest-neighbor similarity nhỏ nhất nhằm giảm thiểu tối đa nguy cơ rò rỉ ngữ nghĩa. Bảng \ref{tab:cegs_split} trình bày cấu hình chi tiết đang được triển khai thực tế.

\begin{table}[htbp]
\centering
\caption{Class-specific configuration of the CEGS-Split protocol}
\label{tab:cegs_split}
\resizebox{\columnwidth}{!}{%
\begin{tabular}{ll}
\toprule
\textbf{Species} & \textbf{Split Strategy} \\
\midrule
\textit{Afzelia africana} & Stratified Group \\
\textit{Afzelia bella} & Hierarchical Clustering \\
\textit{Afzelia pachyloba} & Agglomerative Stratified \\
\textit{Afzelia quanzensis} & Mahalanobis Iterative \\
\textit{Dalbergia cochinchinensis} & Agglomerative Stratified \\
\textit{Dalbergia melanoxylon} & Mahalanobis Iterative \\
\textit{Dalbergia oliveri} & Stratified Group \\
\textit{Dalbergia rimosa} & Hierarchical Clustering \\
\textit{Dalbergia tonkinensis} & Hierarchical Clustering \\
\textit{Guibourtia arnoldiana} & Hierarchical Clustering \\
\textit{Guibourtia coleosperma} & Agglomerative Stratified \\
\textit{Guibourtia ehie} & Hierarchical Clustering \\
\textit{Peltogyne pubescens} & Mahalanobis Iterative \\
\textit{Pterocarpus erinaceus} & Agglomerative Stratified \\
\textit{Pterocarpus indicus} & Agglomerative Stratified \\
\textit{Pterocarpus macrocarpus} & Hierarchical Clustering \\
\textit{Pterocarpus soyauxii} & Hierarchical Clustering \\
\textit{Sindora cochinchinensis} & Mahalanobis Iterative \\
\textit{Sindora tonkinensis} & Agglomerative Stratified \\
\bottomrule
\end{tabular}%
}
\end{table}


Đối với các loài gỗ có phân bổ hoán đổi (Swapped (Val)), kết quả chia tạm thời được điều chỉnh như sau:
\begin{equation}
\mathcal{D}_{\text{train}}=\mathcal{D}_{\text{train\_temp}},\quad
\mathcal{D}_{\text{val}}=\mathcal{D}_{\text{test\_temp}},\quad
\mathcal{D}_{\text{test}}=\mathcal{D}_{\text{val\_temp}}.
\end{equation}
Với các loài có cấu hình tiêu chuẩn (Standard (Test)), kết quả ban đầu được giữ nguyên. Cơ chế này giúp cân bằng độ khó giữa tập kiểm định và tập kiểm thử theo đặc thù của từng loài gỗ, đồng thời vẫn duy trì sự cô lập thư mục con nghiêm ngặt.

\subsection{Feature Backbone and Nonlinear Projection Head}
Mô hình nhận diện chính sử dụng kiến trúc ConvNeXt-Tiny làm mạng trích xuất đặc trưng cơ sở (backbone). Khoảng $90\%$ các lớp đầu tiên được đóng băng (freeze) để giữ lại các đặc trưng tổng quát học từ tập ImageNet, tránh hiện tượng quá khớp (overfitting) trên các loài gỗ có số lượng mẫu vật hạn chế. Đặc trưng thô $\mathbf{f}\in\mathbb{R}^{768}$ sau đó được đưa qua bộ đầu chiếu (projection head) phi tuyến để tạo đặc trưng nhúng $\mathbf{z}\in\mathbb{R}^{256}$. Đối với bộ đầu chiếu gồm hai tầng liên kết đầy đủ, biểu diễn chuẩn hóa $L_2$ được tính toán như sau:
\begin{equation}
\tilde{\mathbf{z}}=W_2\text{ReLU}(W_1\mathbf{f}+b_1)+b_2,
\end{equation}
\begin{equation}
\mathbf{z}=\frac{\tilde{\mathbf{z}}}{\|\tilde{\mathbf{z}}\|_2}.
\end{equation}
Do $\|\mathbf{z}\|_2=1$, khoảng cách Euclid giữa hai đặc trưng nhúng liên hệ trực tiếp với cosine similarity (cosine similarity):
\begin{equation}
d(\mathbf{z}_a,\mathbf{z}_b)=\|\mathbf{z}_a-\mathbf{z}_b\|_2=\sqrt{2(1-\mathbf{z}_a^T\mathbf{z}_b)}.
\end{equation}

\subsection{Deep Metric Learning Loss Functions and Baselines}
Để đánh giá khách quan các hướng học biểu diễn đặc trưng, chúng tôi thực nghiệm đối chiếu 14 hàm loss học khoảng cách trong cùng một quy trình huấn luyện thống nhất về mạng cơ sở, bộ đầu chiếu, tăng cường dữ liệu và giao thức CEGS-Split. Ký hiệu chung: $\mathbf{z}_i\in\mathbb{R}^{D}$ là đặc trưng nhúng đã chuẩn hóa $L_2$, $y_i$ là nhãn loài, $d_{ij}=\|\mathbf{z}_i-\mathbf{z}_j\|_2$ là khoảng cách Euclid, $S_{ij}=\mathbf{z}_i^T\mathbf{z}_j$ là cosine similarity và $[u]_+=\max(0,u)$. Đối với hình ảnh mặt cắt gỗ, các hàm loss này không chỉ được đánh giá bằng accuracy (accuracy) mà còn bằng khả năng xây dựng không gian nhúng có cấu trúc phân cụm ổn định cho các loài gỗ trong cùng chi.

\subsubsection{Standard Triplet Loss}
Triplet Loss tối ưu hóa mối quan hệ giữa bộ ba mẫu gồm mẫu neo (anchor) $a$, mẫu tích cực (positive) $p$ thuộc cùng loài và mẫu tiêu cực (negative) $n$ thuộc loài khác. Mục tiêu của hàm loss là thu hẹp khoảng cách giữa mẫu neo và mẫu tích cực, đồng thời đẩy mẫu tiêu cực ra xa hơn mẫu neo ít nhất một khoảng biên (margin) $\alpha$:
\begin{equation}
L_{\text{triplet}}=\frac{1}{|\mathcal{T}|}\sum_{(a,p,n)\in\mathcal{T}}\left[d(a,p)^2-d(a,n)^2+\alpha\right]_+.
\end{equation}
Đây là mô hình đối chiếu nền tảng để kiểm tra liệu mô hình có học được cấu trúc kéo-gần/đẩy-xa cơ bản hay không. Hạn chế của phương pháp này là số lượng bộ ba hợp lệ tăng nhanh theo kích thước lô dữ liệu $O(B^3)$ và phần lớn các bộ ba trong lô dữ liệu (batch) nhanh chóng trở thành bộ ba dễ (có giá trị loss bằng 0). Trên bộ dữ liệu S3, điều này đặc biệt bất lợi do các loài gỗ cùng chi tạo ra các ranh giới phân tách rất khó, trong khi số lượng mẫu trong một lô dữ liệu có thể không đủ đa dạng để tìm được các bộ ba hữu ích.

\subsubsection{Semi-hard Triplet Loss}
Phương pháp Triplet bán khó (Semi-hard Triplet Loss) lựa chọn các mẫu tiêu cực (negative) không quá dễ nhưng cũng không quá nhiễu. Với một cặp mẫu neo và mẫu tích cực $(a,p)$, mẫu tiêu cực $n$ được chọn nếu thỏa mãn:
\begin{equation}
d(a,p)^2<d(a,n)^2<d(a,p)^2+\alpha.
\end{equation}
Loss vẫn giữ dạng triplet chuẩn:
\begin{equation}
L_{\text{semi-hard}}=\frac{1}{|\mathcal{T}_{sh}|}\sum_{(a,p,n)\in\mathcal{T}_{sh}}\left[d(a,p)^2-d(a,n)^2+\alpha\right]_+.
\end{equation}
Cơ chế này giúp ổn định đạo hàm (gradient) hơn so với việc khai thác mẫu khó nhất (hardest mining), do mẫu tiêu cực vẫn nằm ngoài khoảng cách mẫu tích cực một chút, tránh cho mô hình bị ảnh hưởng bởi các mẫu dị biệt (outlier) hoặc nhãn bị nhiễu. Đối với hình ảnh mặt cắt gỗ, cơ chế khai thác mẫu bán khó (semi-hard mining) rất hữu ích khi cần học biên giới phân biệt giữa các loài gỗ gần nhau trong cùng chi \textit{Dalbergia} mà không làm mô hình bị tối ưu quá mức đối với các mẫu bất thường.

\subsubsection{Soft-Margin Triplet Loss}
Phương pháp Triplet biên mềm (Soft-Margin Triplet) thay thế hàm hinge cứng bằng hàm log-sum-exp mịn:
\begin{equation}
L_{\text{soft-triplet}}=\frac{1}{N}\sum_{i=1}^{N}\log\left[1+e^{d(a_i,p_i)^2-d(a_i,n_i)^2}\right].
\end{equation}
Khác với hàm loss Triplet truyền thống có biên cố định, hàm này vẫn duy trì một lượng cập nhật đạo hàm nhỏ ngay cả khi mẫu tiêu cực đã ở khoảng cách tương đối an toàn. Điều này làm giảm độ nhạy của mô hình đối với việc chọn siêu tham số biên $\alpha$, giúp quá trình tối ưu hội tụ mượt mà hơn. Trong bối cảnh bộ dữ liệu S3, biên mềm đóng vai trò cầu nối hiệu quả giữa phương pháp đối chiếu Triplet cơ bản và các hàm loss hiện đại, giúp mô hình liên tục tinh chỉnh cấu trúc của không gian nhúng thay vì dừng học sớm khi đã vượt qua biên cứng.

\subsubsection{Angular Loss}
Hàm loss Angular (Angular Loss) tối ưu hóa quan hệ góc của tam giác được tạo bởi bộ ba mẫu (triplet) thay vì chỉ tối ưu hóa khoảng cách Euclid đơn thuần. Với bộ ba mẫu $(a,p,n)$, hàm loss được xác định như sau:
\begin{equation}
L_{\text{angular}}=\frac{1}{|\mathcal{T}|}\sum_{(a,p,n)\in\mathcal{T}}\log\left[1+e^{\Delta_{\theta}(a,p,n)}\right],
\end{equation}
where:
\begin{equation}
\begin{aligned}
\Delta_{\theta}(a,p,n) = &4\tan^2\alpha(\mathbf{z}_a+\mathbf{z}_p)^T\mathbf{z}_n \\
&- 2(1+\tan^2\alpha)\mathbf{z}_a^T\mathbf{z}_p.
\end{aligned}
\end{equation}
Do đặc trưng nhúng đã được chuẩn hóa L2, biên góc ít bị ảnh hưởng bởi độ lớn (scale) của vector đặc trưng. Điều này rất phù hợp với hình ảnh gỗ vĩ mô, nơi cùng một cấu trúc thớ gỗ có thể thay đổi cường độ sáng, độ phóng đại hoặc độ ẩm bề mặt, nhưng quan hệ góc trong không gian nhúng vẫn bảo toàn được thông tin phân biệt loài.

\subsubsection{Multi-Similarity Loss}
Multi-Similarity Loss khai thác đồng thời thông tin từ nhiều cặp tích cực và tiêu cực trong cùng một lô dữ liệu, đồng thời gán trọng số lớn hơn cho các cặp có độ khó cao:
\begin{equation}
L_{\text{MS}}=\frac{1}{N}\sum_{i=1}^{N} \left[ \phi^+(P_i) + \phi^-(N_i) \right],
\end{equation}
where:
\begin{equation}
\begin{aligned}
\phi^+(P_i) &= \frac{1}{\alpha}\log\left[1+\sum_{p\in P_i}e^{-\alpha(S_{ip}-\lambda)}\right], \\
\phi^-(N_i) &= \frac{1}{\beta}\log\left[1+\sum_{n\in N_i}e^{\beta(S_{in}-\lambda)}\right].
\end{aligned}
\end{equation}
Thành phần tích cực phạt các mẫu cùng loài chưa đủ gần nhau trong không gian đặc trưng; thành phần tiêu cực phạt các mẫu khác loài nhưng lại có cosine similarity lớn. Đây là một trong những hàm loss đối chiếu quan trọng nhất đối với bộ dữ liệu S3, do các loài gỗ khác nhau nhưng cùng chi thường tạo ra các mẫu tiêu cực khó (hard negative) tự nhiên trong thực tế, ví dụ giữa loài \textit{Dalbergia oliveri} và \textit{Dalbergia tonkinensis}.

\subsubsection{Standard Contrastive Loss}
Contrastive Loss tối ưu hóa khoảng cách trên từng cặp mẫu riêng lẻ. Gọi $q_{ij}=0$ nếu hai mẫu $i, j$ cùng loài ($y_i=y_j$) và $q_{ij}=1$ nếu khác loài ($y_i\ne y_j$), hàm loss được phát biểu:
\begin{equation}
L_{\text{contrastive}}=(1-q_{ij})\frac{1}{2}d_{ij}^{2}+q_{ij}\frac{1}{2}[m-d_{ij}]_+^2.
\end{equation}
Hàm loss kéo các cặp hình ảnh cùng loài về gần nhau và đẩy các cặp khác loài ra ngoài khoảng biên $m$. Ưu điểm của phương pháp này là trực quan và dễ cài đặt; tuy nhiên, hạn chế lớn là một biên khoảng cách duy nhất được áp dụng chung cho mọi cặp mẫu tiêu cực, không phản ánh được mức độ khó phân biệt khác nhau giữa các loài cùng chi so với các loài thuộc các chi khác nhau.

\subsubsection{Online Hard Contrastive Loss}
Phương pháp tương phản khai thác mẫu khó trực tuyến (Online Hard Contrastive) tập trung tối ưu hóa trên cặp tích cực xa nhất và cặp tiêu cực gần nhất trong cùng một lô dữ liệu (batch):
\begin{equation}
d_{\text{pos\_max}}^i=\max_{p:y_p=y_i,p\ne i}d_{ip},\quad d_{\text{neg\_min}}^i=\min_{n:y_n\ne y_i}d_{in}.
\end{equation}
\begin{equation}
L_{\text{hard-cont}}=\frac{1}{B}\sum_{i=1}^{B}\left[(d_{\text{pos\_max}}^i)^2+[m-d_{\text{neg\_min}}^i]_+^2\right].
\end{equation}
Cách tiếp cận này tập trung tài nguyên tính toán đạo hàm vào các vùng ranh giới khó phân tách nhất thay vì các cặp mẫu dễ. Đối với hình ảnh gỗ, nó giúp mô hình tập trung phân biệt các mẫu gỗ có vân tương đồng giữa hai loài khác nhau, nhưng cũng nhạy cảm hơn với các mẫu dị biệt (outlier), ảnh bị nhiễu hoặc nhãn bị gán sai.

\subsubsection{Lifted Structured Loss}
Lifted Structured Loss mở rộng cơ chế học tương phản bằng cách đồng thời xem xét tất cả các mẫu tiêu cực liên quan đến một cặp mẫu tích cực $(i,j)$:
\begin{equation}
L_{\text{lifted}}=\frac{1}{2|P|}\sum_{(i,j)\in P}\left[d_{ij}+\log\left(\Phi(i, j)\right)\right]_+^2,
\end{equation}
where:
\begin{equation}
\begin{aligned}
\Phi(i, j) = &\sum_{k:y_k\ne y_i}e^{\alpha-d_{ik}} \\
&+ \sum_{k:y_k\ne y_j}e^{\alpha-d_{jk}}.
\end{aligned}
\end{equation}
Hàm $\log\sum\exp$ đóng vai trò là một xấp xỉ trơn của phép toán lấy giá trị cực đại, do đó các mẫu tiêu cực gần mẫu neo nhất sẽ tự động nhận lực đẩy mạnh nhất. Hàm loss này rất hữu ích để đánh giá hiệu quả của việc tối ưu hóa cấu trúc lô dữ liệu đối với bề mặt gỗ, nơi nhiều loài khác nhau thường có một số mẫu đặc trưng nằm rất gần nhau trong không gian nhúng.

\subsubsection{Circle Loss}
Circle Loss đề xuất một cách tiếp cận tối ưu hóa độ tương đồng thống nhất, trong đó các mẫu tích cực và tiêu cực được gán các trọng số tự điều chỉnh:
\begin{equation}
L_{\text{circle}} = \log \left[ 1 + \Omega_{\text{circle}}(i) \right],
\end{equation}
where:
\begin{equation}
\begin{aligned}
\Omega_{\text{circle}}(i) = &\sum_{n\in N_i}e^{\gamma\alpha_n(S_{in}-\Delta_n)} \times \\
&\sum_{p\in P_i}e^{-\gamma\alpha_p(S_{ip}-\Delta_p)},
\end{aligned}
\end{equation}
and:
\begin{equation}
\alpha_p=[O_p-S_{ip}]_+, \quad \alpha_n=[S_{in}-O_n]_+.
\end{equation}
Mẫu tích cực đã ở đủ gần sẽ nhận trọng số cập nhật nhỏ, trong khi mẫu tiêu cực quá gần sẽ bị đẩy đi rất mạnh. Cơ chế này đặc biệt phù hợp với bài toán phân loại chi tiết (fine-grained), do nó tránh việc tối ưu hóa quá mức trên các loài dễ phân biệt, đồng thời tập trung đạo hàm vào các cặp mẫu thuộc cùng một chi có độ tương đồng cao.

\subsubsection{Proxy Anchor Loss}
Phương pháp Proxy Anchor thay thế việc so sánh trực tiếp giữa các hình ảnh bằng việc so sánh hình ảnh với các mẫu đại diện (proxy) của từng lớp. Với proxy $p$ đại diện cho một lớp, hàm loss được tính như sau:
\begin{equation}
L_{\text{PA}} = \frac{1}{|P^+|}\sum_{p\in P^+}\log \Gamma^+(p) + \frac{1}{|P|}\sum_{p\in P}\log \Gamma^-(p),
\end{equation}
where:
\begin{equation}
\begin{aligned}
\Gamma^+(p) &= 1+\sum_{\mathbf{z}\in X_p^+}e^{-\alpha(S(\mathbf{z},p)-\delta)}, \\
\Gamma^-(p) &= 1+\sum_{\mathbf{z}\in X_p^-}e^{\alpha(S(\mathbf{z},p)+\delta)}.
\end{aligned}
\end{equation}
Do số lượng lớp của bộ dữ liệu S3 tương đối nhỏ, phương pháp học dựa trên proxy (proxy-based learning) mang lại chi phí tính toán thấp và dòng đạo hàm ổn định. Ngoài ra, Proxy Anchor giúp giảm thiểu sự phụ thuộc vào cấu trúc thành phần của lô dữ liệu (batch composition) - một ưu điểm quan trọng khi các loài gỗ trong cơ sở dữ liệu có sự mất cân bằng lớn về số lượng ảnh và số lượng thư mục con.

\subsubsection{ArcFace Loss}
Phương pháp ArcFace (Additive Angular Margin Loss) áp dụng trực tiếp một biên góc cộng (additive angular margin) vào giá trị logit của lớp chính xác:
\begin{equation}
\begin{aligned}
L_{\text{ArcFace}} = -\frac{1}{N}\sum_{i=1}^{N}\log \frac{e^{s\cos(\theta_{y_i}+m)}}{\Lambda_i},
\end{aligned}
\end{equation}
where:
\begin{equation}
\begin{aligned}
\Lambda_i = e^{s\cos(\theta_{y_i}+m)} + \sum_{j\ne y_i}e^{s\cos\theta_j}.
\end{aligned}
\end{equation}
Do cả đặc trưng nhúng và ma trận trọng số của lớp đều được chuẩn hóa, ArcFace thiết lập một biên khoảng cách trắc địa (geodesic margin) trên bề mặt siêu cầu đơn vị. Hàm loss này tỏ ra đặc biệt hiệu quả trong bài toán phân loại chi tiết (fine-grained), do nó ép cụm đặc trưng của cùng một lớp co lại trong một vùng góc hẹp và đẩy cụm của các lớp khác ra xa, rất thích hợp với các loài gỗ có sự khác biệt hình thái cực kỳ nhỏ.

\subsubsection{SubCenter ArcFace Loss}
Phương pháp SubCenter ArcFace mở rộng thuật toán ArcFace bằng cách gán nhiều trọng tâm phụ (sub-center) cho mỗi lớp. Với cấu hình $K$ trọng tâm phụ cho mỗi loài, giá trị logit của lớp chính xác được tính toán dựa trên trọng tâm phụ gần nhất:
\begin{equation}
\cos\theta_{y_i}=\max_{k=1}^{K}\frac{\mathbf{z}_i^T W_{y_i}^{(k)}}{\|\mathbf{z}_i\|_2\|W_{y_i}^{(k)}\|_2}.
\end{equation}
Sau đó, mô hình áp dụng cùng một biên góc tương tự như ArcFace. Cơ chế đa tâm này rất phù hợp với đặc thù bộ dữ liệu S3, nơi các hình ảnh của một loài gỗ có thể tạo ra phân bố đa đỉnh trong không gian đặc trưng do sự khác biệt giữa gỗ lõi và gỗ giác, tuổi cây hoặc vị trí cắt. Thay vì ép buộc tất cả các biến thể tự nhiên này vào một trọng tâm duy nhất, SubCenter ArcFace cho phép mô hình tự động học các cụm phân bố cấu trúc bề mặt riêng biệt trong cùng một loài gỗ.

\subsubsection{SoftTriple Loss}
Tương tự, phương pháp SoftTriple cũng sử dụng nhiều trọng tâm đại diện cho mỗi lớp nhưng áp dụng cơ chế gán mềm (soft assignment) thay vì lấy giá trị cực đại:
\begin{equation}
S_{i,c}=\sum_{k=1}^{K}\frac{e^{\mathbf{z}_i^T w_c^{(k)}/\gamma}}{\sum_{t=1}^{K}e^{\mathbf{z}_i^T w_c^{(t)}/\gamma}}\mathbf{z}_i^T w_c^{(k)}.
\end{equation}
Loss phân loại có margin:
\begin{equation}
L_{\text{SoftTriple}}=-\log\frac{e^{\lambda(S_{i,y_i}-\delta)}}{e^{\lambda(S_{i,y_i}-\delta)}+\sum_{c\ne y_i}e^{\lambda S_{i,c}}}.
\end{equation}
Cơ chế gán mềm giúp dòng đạo hàm cập nhật đồng thời cho nhiều trọng tâm trong mỗi bước huấn luyện, hạn chế hiện tượng trọng tâm bị cô lập hoặc không được cập nhật (dead center). Đối với hình ảnh gỗ, SoftTriple mang lại ý nghĩa thực tế lớn do các phân nhóm vân gỗ thường không phân tách một cách rõ rệt mà có sự chuyển tiếp liên tục giữa các vùng mô trên mặt cắt.

\subsubsection{Supervised Contrastive Loss (SupCon)}
Contrastive Loss có giám sát (Supervised Contrastive Loss - SupCon) mở rộng phương pháp SimCLR sang học có giám sát bằng cách coi tất cả các mẫu thuộc cùng một loài trong lô dữ liệu là các mẫu tích cực của nhau:
\begin{equation}
L_{\text{SupCon}}=\sum_{i\in I}\frac{-1}{|P(i)|}\sum_{p\in P(i)}\log\frac{e^{\mathbf{z}_i^T\mathbf{z}_p/\tau}}{\sum_{a\in A(i)}e^{\mathbf{z}_i^T\mathbf{z}_a/\tau}}.
\end{equation}
SupCon tối ưu hóa việc sắp đặt đặc trưng trên toàn bộ lô dữ liệu và thường tạo ra không gian nhúng có khả năng chuyển giao tốt. Trong bài toán S3, phương pháp này giúp kéo hình ảnh của cùng một loài lại gần nhau bất kể những biến thiên do góc chụp, nhưng vẫn tiềm ẩn nguy cơ đẩy các mẫu có độ tương đồng cao sang hai phía quá mạnh (false-negative repulsion) when two species under the same genus are extremely similar.

\subsection{Self-Supervised Learning Loss Functions}
Trong học biểu diễn tự giám sát (Self-Supervised Learning - SSL), mô hình thực hiện tối ưu hóa đặc trưng thông qua các biến thể tăng cường ngẫu nhiên của hình ảnh mà không cần sử dụng nhãn lớp giám sát. Phần này trình bày các hàm loss của bốn phương pháp học tự giám sát được khảo sát trong nghiên cứu.

\subsubsection{SimCLR Loss}
Phương pháp SimCLR (Simple Framework for Contrastive Learning of Visual Representations) tối ưu hóa hàm loss tương phản InfoNCE trên các cặp ảnh tăng cường. Với kích thước lô dữ liệu là $N$, mỗi hình ảnh được tăng cường ngẫu nhiên thành 2 phiên bản (views), tạo thành một lô gồm $2N$ mẫu dữ liệu. Cặp mẫu tích cực (positive pair) là $(i,j)$ gồm hai phiên bản tăng cường khác nhau của cùng một hình ảnh gốc. Hàm loss cho mỗi cặp mẫu tích cực $(i,j)$ được phát biểu:
\begin{equation}
\begin{aligned}
\ell(i, j) = -\log \frac{e^{\mathbf{z}_i^T\mathbf{z}_j/\tau}}{\Omega_i},
\end{aligned}
\end{equation}
where:
\begin{equation}
\begin{aligned}
\Omega_i = \sum_{k=1}^{2N} \mathbb{I}_{[k \ne i]} e^{\mathbf{z}_i^T\mathbf{z}_k/\tau}.
\end{aligned}
\end{equation}
trong đó $\tau$ là siêu tham số nhiệt độ (temperature). Hàm loss tổng thể được tính bằng macro average của toàn bộ các cặp mẫu tích cực trong lô dữ liệu:
\begin{equation}
L_{\text{SimCLR}} = \frac{1}{2N} \sum_{k=1}^{N} \left[ \ell(2k-1, 2k) + \ell(2k, 2k-1) \right].
\end{equation}

\subsubsection{BYOL Loss}
Phương pháp BYOL (Bootstrap Your Own Latent) sử dụng hai mạng nơ-ron song song: mạng học trực tiếp (online network) với tham số $\theta$ và target network với tham số $\xi$. Mạng học trực tiếp có nhiệm vụ dự đoán biểu diễn của mạng mục tiêu từ một phiên bản ảnh khác. BYOL tối ưu hóa sai số bình phương trung bình (MSE) giữa vector dự đoán đã chuẩn hóa của mạng trực tiếp và vector đặc trưng của mạng mục tiêu:
\begin{equation}
L_{\text{MSE}}(q_\theta(z_\theta), z'_\xi) = \left\| \bar{q}_\theta - \bar{z}'_\xi \right\|_2^2 = 2 - 2 \bar{q}_\theta^T \bar{z}'_\xi,
\end{equation}
where:
\begin{equation}
\bar{q}_\theta = \frac{q_\theta(z_\theta)}{\|q_\theta(z_\theta)\|_2},\quad \bar{z}'_\xi = \frac{z'_\xi}{\|z'_\xi\|_2}.
\end{equation}
Việc ngăn chặn hiện tượng sụp đổ biểu diễn đặc trưng (representation collapse) được đảm bảo bằng cơ chế dừng cập nhật đạo hàm (\textit{stop-gradient}) đặt trên mạng mục tiêu:
\begin{equation}
L_{\text{BYOL}} = L_{\text{MSE}}(q_\theta(z_\theta), \text{stop\_gradient}(z'_\xi)).
\end{equation}
Tham số mạng mục tiêu $\xi$ được cập nhật chậm theo cơ chế trung bình trượt lũy thừa (EMA) từ $\theta$:
\begin{equation}
\xi \leftarrow m \xi + (1-m) \theta,
\end{equation}
trong đó $m \in [0, 1]$ là hệ số momentum.

\subsubsection{SimSiam Loss}
Phương pháp SimSiam (Simple Siamese) đơn giản hóa kiến trúc BYOL bằng cách loại bỏ hoàn toàn momentum target network (momentum target network, tức thiết lập $\xi = \theta$). SimSiam chứng minh rằng chỉ cần sử dụng cơ chế dừng cập nhật đạo hàm kết hợp với kiến trúc không đối xứng (bộ dự đoán predictor) là đủ để ngăn chặn sự sụp đổ biểu diễn đặc trưng. Hàm loss đối xứng được định nghĩa như sau:
\begin{equation}
L_{\text{SimSiam}} = \frac{1}{2} \mathcal{D}\big(p_1, \text{sg}(z_2)\big) + \frac{1}{2} \mathcal{D}\big(p_2, \text{sg}(z_1)\big),
\end{equation}
trong đó $p_1 = q_\theta(z_1)$ và $p_2 = q_\theta(z_2)$ là đầu ra của bộ dự đoán cho hai phiên bản ảnh, còn $z_1, z_2$ là đầu ra của bộ đầu chiếu. $\text{sg}$ biểu thị toán tử dừng đạo hàm (stop-gradient). $\mathcal{D}$ là cosine similarity âm:
\begin{equation}
\mathcal{D}(p, z) = -\frac{p^T z}{\|p\|_2 \|z\|_2}.
\end{equation}

\subsubsection{Barlow Twins Loss}
Phương pháp Barlow Twins áp dụng nguyên lý giảm thiểu thông tin dư thừa (redundancy reduction) lên ma trận tương quan chéo thu được từ hai phiên bản ảnh tăng cường. Gọi $C$ là ma trận tương quan chéo kích thước $D \times D$ tính toán dọc theo chiều lô dữ liệu $B$ giữa hai tập đặc trưng chiếu $\mathbf{z}^A$ và $\mathbf{z}^B$:
\begin{equation}
C_{ij} = \frac{\sum_{b=1}^{B} z_{b,i}^A z_{b,j}^B}{\sqrt{\sum_{b=1}^{B} (z_{b,i}^A)^2} \sqrt{\sum_{b=1}^{B} (z_{b,j}^B)^2}},
\end{equation}
trong đó $i, j$ đại diện cho các chiều của không gian đặc trưng nhúng. Hàm mất mát Barlow Twins được định nghĩa là:
\begin{equation}
L_{\text{BarlowTwins}} = \sum_{i=1}^{D} (1 - C_{ii})^2 + \lambda \sum_{i=1}^{D} \sum_{j \neq i} C_{ij}^2,
\end{equation}
trong đó số hạng thứ nhất (invariance term) ép các hệ số tương quan chéo trên đường chéo chính tiến về 1 (tối đa hóa độ tương đồng giữa hai phiên bản ảnh), số hạng thứ hai (redundancy reduction term) ép các hệ số tương quan chéo ngoài đường chéo chính tiến về 0 với hệ số phạt $\lambda > 0$ (giảm thiểu sự trùng lặp thông tin giữa các chiều đặc trưng nhúng).

